We need data. But not just any data --- we need data we can \emph{trust}
to be fair, and data we can carve up to test the model honestly.

A teacher preparing worksheets for a test is the model. A bad teacher
writes twenty questions that are all the same question with different
numbers --- the student can memorize one template. A good teacher writes
questions that look different on the surface but test the same idea, so
a student who actually learned the idea passes whether or not they've
seen the phrasing.

Our \textbf{teacher} is that good worksheet-writer. But there's
something special about the way we build it, and it's the single most
important design decision in the whole book.

\section{Semantic-first: start with the meaning, not a
template}\label{semantic-first-start-with-the-meaning-not-a-template}

The naive way to generate data is to write a template and fill in a few
blanks:

\begin{verbatim}
template -> sentence -> code
\end{verbatim}

We write ``add \{a\} and \{b\}'' and a few code snippets, and we are
done. This is fast and easy, and it's exactly the wrong way to do it.
The problem is that the sentence and the code are generated
\textbf{independently}, from a template, so there's no guarantee they
mean the same thing. We can end up with a sentence that says ``add'' but
code that subtracts, or a test that checks the wrong thing.

We do the opposite. We start from the \textbf{meaning} and generate
everything else \textbf{from it}:

\begin{verbatim}
SEMANTIC INTENT
   -> PROGRAM STRUCTURE
      -> REFERENCE CODE
      -> NATURAL LANGUAGE
      -> HIDDEN TESTS
\end{verbatim}

The key insight is at the top: \textbf{one semantic record drives all
three.} There is no way for the sentence to disagree with the code or
the tests, because they are all the same meaning, rendered three ways.
If we start with ``the function should return \texttt{a\ +\ b},'' the
sentence says ``sum,'' the code adds, and the tests check addition ---
and they can't drift apart, because they all came from the same seed.

\section{What a ``semantic record'' looks
like}\label{what-a-semantic-record-looks-like}

The meaning is captured in a small, structured object. It is not English
and not code; it is the \emph{idea} underneath both. For an addition
task it looks like this:

\begin{Shaded}
\begin{Highlighting}[]
\FunctionTok{\{}
  \DataTypeTok{"operation"}\FunctionTok{:} \StringTok{"add"}\FunctionTok{,}
  \DataTypeTok{"args"}\FunctionTok{:} \OtherTok{[}\FunctionTok{\{}\DataTypeTok{"name"}\FunctionTok{:} \StringTok{"a"}\FunctionTok{,} \DataTypeTok{"type"}\FunctionTok{:} \StringTok{"int"}\FunctionTok{,} \DataTypeTok{"role"}\FunctionTok{:} \StringTok{"operand"}\FunctionTok{\}}\OtherTok{,}
           \FunctionTok{\{}\DataTypeTok{"name"}\FunctionTok{:} \StringTok{"b"}\FunctionTok{,} \DataTypeTok{"type"}\FunctionTok{:} \StringTok{"int"}\FunctionTok{,} \DataTypeTok{"role"}\FunctionTok{:} \StringTok{"operand"}\FunctionTok{\}}\OtherTok{]}\FunctionTok{,}
  \DataTypeTok{"output\_type"}\FunctionTok{:} \StringTok{"int"}
\FunctionTok{\}}
\end{Highlighting}
\end{Shaded}

From this single object, we generate:

\begin{enumerate}
\def\labelenumi{\arabic{enumi}.}
\tightlist
\item
  \textbf{the reference program} ---
  \texttt{def\ solve(a,\ b):\ return\ a\ +\ b},
\item
  \textbf{the natural-language prompt} --- e.g.~\emph{``Write a function
  that takes two numbers, a and b, and returns their sum''},
\item
  \textbf{the hidden tests} ---
  e.g.~\texttt{\{args:{[}-28,47{]},\ expected:19\}},
  \texttt{\{args:{[}0,0{]},\ expected:0\}}, and edge cases we'll get to.
\end{enumerate}

This is the heart of the teacher. The pump is \texttt{make\_task}, which
takes an operation and a seed, builds the semantics, and then calls
three sub-generators that all read that one record.

\section{The four sub-generators}\label{the-four-sub-generators}

The teacher is really four small pieces working together, all fed by the
same semantics:

\textbf{1. \texttt{make\_semantics}} --- builds the semantic intent: the
operation plus its argument descriptors. It pulls argument names from a
controlled set (so we can vary them later), assigns an input/output
type, and tags the task with a complexity score.

\textbf{2. \texttt{generate\_reference}} --- turns the semantics into
the ground-truth Python. This is the correct program the grader is
validated against. For \texttt{add} it emits
\texttt{def\ solve(a,b):\ return\ a+b}; for \texttt{list\_sum} it emits
a loop over the list.

\textbf{3. \texttt{render\_nl}} --- turns the same semantics into
English, using one of several \textbf{template families}. A template
family is just a different way of phrasing the same instruction. For
addition, we have phrases like:

\begin{itemize}
\tightlist
\item
  \emph{``Write a function that takes two numbers, a and b, and returns
  their sum.''}
\item
  \emph{``Implement solve so that given a and b it returns their sum.''}
\item
  \emph{``Define a function solve(a, b) that returns a + b.''}
\end{itemize}

Same meaning, three different wordings. This is what gives the dataset
linguistic diversity --- and, more importantly, it's what lets us later
test whether the model learned the \emph{idea} or just one
\emph{sentence}.

\textbf{4. \texttt{generate\_hidden\_tests}} --- generates test cases
\textbf{from the semantics}, not from the reference program's logic. It
deliberately covers edge cases: ordinary values, zero, negatives,
boundary values, large magnitudes, and (where legal) empty lists. For
each it computes the expected value from the operation itself.

\section{Why the tests come from the meaning, not the
code}\label{why-the-tests-come-from-the-meaning-not-the-code}

This is a subtle but crucial point. If we generated the tests by looking
at the reference program, a model that memorized the reference would
already know what the tests check. By generating them from the
\emph{semantics} instead, the tests probe whether the model
independently understands the operation --- whether it can produce the
right answer for inputs it could not have seen in the reference.

\section{Maximizing controlled
diversity}\label{maximizing-controlled-diversity}

A good teacher does not just produce many examples; it produces
\emph{diverse} ones while holding the semantics fixed. We deliberately
vary three things:

\begin{itemize}
\tightlist
\item
  \textbf{constants} --- different values each task, so the model can't
  memorize outputs,
\item
  \textbf{identifiers} --- \texttt{a,\ b} vs \texttt{left,\ right}, so
  it can't tie the operation to particular names,
\item
  \textbf{wording} --- several template families per operation, so it
  sees the same idea said many ways.
\end{itemize}

This controlled diversity is what makes the \emph{generalization tests}
later possible. Because we control \emph{how} we vary the surface, we
can also \emph{hold it out} and test whether the model generalizes.
That's not an accident; it's the whole point.

\section{Avoiding accidental
correlation}\label{avoiding-accidental-correlation}

The danger in any generated dataset is \textbf{accidental correlation}
--- a pattern that makes the answer trivially predictable from the
surface. The classic failures:

\begin{itemize}
\tightlist
\item
  every addition example uses the word ``total'';
\item
  every maximum example uses the variable \texttt{x};
\item
  every multiplication has exactly three arguments;
\item
  every loop example is exactly one sentence long;
\item
  every operation uses one unique wording.
\end{itemize}

If those correlations exist, a model (or a trivial classifier) can solve
the task without learning the computation. The generator is built to
break them, and then --- this is the crucial part --- we \textbf{check
that it actually did}.

\section{The leakage audit}\label{the-leakage-audit}

Before trusting any dataset, we ask: \emph{can a trivial classifier
predict the answer from the surface alone?} The method is deliberately
crude: train a TF-IDF + logistic regression model on the
natural-language prompts and see how well it predicts the target
operation from the words alone. If a dumb classifier succeeds, the
wording is leaking the answer.

And it caught us. On our first pilot renderer, a trivial classifier
predicted the operation from the wording with \textbf{100\% accuracy}.
That's a red flag, and it's a valuable catch: it means our wording was
so tied to the operation that a model which ``learned'' might just be
exploiting the phrasing. We recorded it as a findable artifact
(finding-001) and noted exactly what it means for the experiments that
follow.

This is the teacher doing its real job. A dataset generator that only
produces data is half a tool. One that produces data \textbf{and audits
it} is a laboratory instrument.

\section{The takeaway}\label{the-takeaway}

The teacher is not a bulk-data tool; it is a controlled instrument. Its
job is to produce data with \emph{controlled} diversity and
\emph{audited} cleanliness, so that when the model succeeds we can
attribute it to learning the computation --- not to the data
accidentally giving away the answer.

And the reason this matters so much is the whole book's premise.
Everything later --- the generalization tests, the composition test, the
capability map --- depends on the data being a \emph{fair} test. The
teacher is where that fairness is built.
